\section{The main convergence theorem}
\label{sec:main_theorem}

We now state and prove the test-time scaling law: the probability that
the layer-$i$, head-$h$ attention output $x_T$ at the
$\langle/\mathrm{think}\rangle$ position fails to land in the
decoding-margin ball around $V^*(Q)$ decays exponentially in the
reasoning horizon $T$, with rate $p_0 / 8$ set by the anchor-emission
probability of the policy.

\begin{theorem}[Test-time scaling for anchored attention]\label{thm:main_convergence_hp}
Suppose
\Cref{ass:anchor_set_accuracy,ass:anchor_emission_prob,ass:score_margin,ass:bounded_value_norms,ass:decoding_existence}
hold, and assume the score margin $\Delta$ of \Cref{ass:score_margin}
satisfies
\begin{align}\label{eq:thm_main_Delta_condition}
   \Delta
   \;\ge\;
   \log\!\frac{4(M + \max_{Q \in F}\norm{V^*(Q)})}{p_0\, \gamma_{\min}}.
\end{align}
(Equivalently, the score margin must be large enough to drive the
non-anchor leakage below $\gamma_{\min}/2$.) Then for every $Q \in F$
and every reasoning horizon $T \ge 1$,
\begin{align}\label{eq:test_time_scaling}
   \Pr\!\Big[\, \norm{x_T - V^*(Q)} \;\le\; \tfrac{\gamma(Q)}{2}
                + \varepsilon_{\mathrm{anc}}
              \;\;\text{and}\;\;
              \dec(x_T) \in \Correct(Q) \,\Big]
   \;\ge\;
   1 - 2 \exp\!\Big( -\frac{p_0 T}{8} \Big).
\end{align}
In particular, the failure probability decays exponentially in the
test-time horizon $T$:
$\Pr[\text{failure}] \le 2 \exp(-p_0 T / 8) \to 0$ as $T \to \infty$.
\end{theorem}

\begin{proof}
The proof has three steps: discharge the score-margin hypothesis,
identify the single probabilistic event the bound depends on, and read
off the rate. We may assume $2 \exp(-p_0 T/8) < 1$ (else the bound is
vacuous), so $\delta \coloneqq 2 \exp(-p_0 T/8) \in (0, 1)$ is a
valid confidence parameter.

\noindent\textbf{Step 1: deterministic-on-event bound via
\Cref{lem:T_polynomial}.}
\Cref{eq:thm_main_Delta_condition} implies the per-question score-margin
condition \Cref{eq:Delta_condition} for every $Q \in F$, because
$\gamma(Q) \ge \gamma_{\min}$ and
$\norm{V^*(Q)} \le \max_{Q \in F} \norm{V^*(Q)}$ (so the right-hand side
of \Cref{eq:thm_main_Delta_condition} dominates the right-hand side of
\Cref{eq:Delta_condition}). With $\delta = 2 \exp(-p_0 T/8)$ as above,
$\log(2/\delta) = \log(\exp(p_0 T / 8)) = p_0 T / 8$ by
construction, so the lemma's horizon
$T(Q, \delta) = (8/p_0) \log(2/\delta) = (8/p_0) \cdot (p_0 T/8) = T$,
and the hypothesis $T \ge T(Q, \delta)$ of \Cref{lem:T_polynomial} holds
tautologically. The lemma then gives
\begin{align}\label{eq:thm_main_T_poly_conclusion}
   \Pr\!\Big[\, \norm{x_T - V^*(Q)} \;\le\; \tfrac{\gamma(Q)}{2}
              \;+\; \varepsilon_{\mathrm{anc}} \,\Big]
   \;\ge\; 1 - \tfrac{\delta}{2}
   \;=\; 1 - \exp\!\Big( -\frac{p_0 T}{8} \Big).
\end{align}
On the event $\{\norm{x_T - V^*(Q)} \le \gamma(Q)/2 + \varepsilon_{\mathrm{anc}}\}$,
\Cref{lem:decoding_correctness} gives both
$\norm{x_T - V^*(Q)} \le \gamma(Q)$ and $\dec(x_T) \in \Correct(Q)$, so
this event is contained in the joint event of
\Cref{eq:test_time_scaling}.

\noindent\textbf{Step 2: union-bound paragraph.}
The only probabilistic event used in Step 1 is the anchor-count event
$\Ecal_1 \coloneqq \{|\mathcal A^{\mathrm{traj}}_T| \ge p_0 T / 2\}$ that
underlies the proof of \Cref{lem:T_polynomial} via
\Cref{lem:anchor_count_lb}. By a union bound over $\Ecal_1$ alone (the
only random event with non-zero failure probability), the failure
probability of the joint event of \Cref{eq:test_time_scaling} is at
most
\begin{align}\label{eq:thm_main_failure_budget}
   \Pr[\Ecal_1^c]
   \;\le\;
   \tfrac{\delta}{2}
   \;=\;
   \exp\!\Big( -\frac{p_0 T}{8} \Big)
   \;\le\;
   2 \exp\!\Big( -\frac{p_0 T}{8} \Big).
\end{align}

\noindent\textbf{Step 3: conclude.}
Combining Steps 1 and 2, the joint event of \Cref{eq:test_time_scaling}
has probability at least
$1 - \Pr[\Ecal_1^c] \ge 1 - 2 \exp(-p_0 T / 8)$, as desired.
(The factor $2$ is slack absorbed for compatibility with
\Cref{cor:expected_error_scaling}; the tighter bound
$1 - \exp(-p_0 T/8)$ also holds.) When
$2 \exp(-p_0 T/8) \ge 1$, the conclusion holds vacuously since the
right-hand side of \Cref{eq:test_time_scaling} is non-positive.
\end{proof}

\begin{corollary}[Expected-error decay]\label{cor:expected_error_scaling}
Under the assumptions of \Cref{thm:main_convergence_hp}, for every
$Q \in F$ and every $T \ge 1$,
\begin{align}\label{eq:expected_error_scaling}
   \E\!\bigl[\, \norm{x_T - V^*(Q)} \,\bigr]
   \;\le\;
   \tfrac{\gamma(Q)}{2} + \varepsilon_{\mathrm{anc}}
   \;+\;
   2 \bigl(M + \norm{V^*(Q)}\bigr) \exp\!\Big( -\frac{p_0 T}{8} \Big).
\end{align}
The right-hand side has a $T$-independent floor
$\gamma(Q)/2 + \varepsilon_{\mathrm{anc}}$ set by the trained-model
quality, plus a term that decays exponentially in the test-time horizon
$T$. As $T \to \infty$,
$\E[\norm{x_T - V^*(Q)}] \to \gamma(Q)/2 + \varepsilon_{\mathrm{anc}}$.
\end{corollary}

\begin{proof}
The proof is the standard tail-to-expectation conversion. Let
$Y \coloneqq \norm{x_T - V^*(Q)}$ and let $G \coloneqq M + \norm{V^*(Q)}$.
By the triangle inequality and \Cref{ass:bounded_value_norms},
$Y \le \norm{x_T} + \norm{V^*(Q)} \le M + \norm{V^*(Q)} = G$ almost
surely: indeed, by \Cref{lem:softmax_running_average},
$x_T = \sum_k w_{T,k} V_k$ with $\sum_k w_{T,k} = 1$ and $w_{T,k} \ge 0$,
so $x_T$ is a convex combination of the value vectors $\{V_k\}$, each
with $\norm{V_k} \le M$ (by \Cref{ass:bounded_value_norms}); triangle
inequality on this convex combination gives
$\norm{x_T} \le \sum_k w_{T,k} \norm{V_k} \le M$.

Let $\Ecal^\star$ denote the success event of
\Cref{eq:test_time_scaling}: on $\Ecal^\star$,
$Y \le \gamma(Q)/2 + \varepsilon_{\mathrm{anc}}$; on the complement,
$Y \le G$ deterministically. Then
\begin{align*}
   \E[Y]
   \;=\; \E[Y \1\{\Ecal^\star\}] + \E[Y \1\{(\Ecal^\star)^c\}]
   \;\le\; \bigl(\tfrac{\gamma(Q)}{2} + \varepsilon_{\mathrm{anc}}\bigr)
           \cdot \Pr[\Ecal^\star]
         + G \cdot \Pr[(\Ecal^\star)^c]
   \;\le\;
   \tfrac{\gamma(Q)}{2} + \varepsilon_{\mathrm{anc}}
         + G \cdot \Pr[(\Ecal^\star)^c],
\end{align*}
where the first equality is the law of total expectation, the first
inequality uses $Y \le \gamma(Q)/2 + \varepsilon_{\mathrm{anc}}$ on
$\Ecal^\star$ and $Y \le G$ on $(\Ecal^\star)^c$, and the last
inequality drops the $\Pr[\Ecal^\star] \le 1$ multiplier on the success
term. Applying \Cref{thm:main_convergence_hp}
($\Pr[(\Ecal^\star)^c] \le 2 \exp(-p_0 T/8)$) gives
\Cref{eq:expected_error_scaling}, as desired.
\end{proof}

\begin{corollary}[Next-token entropy decay]\label{cor:entropy_decay}
Suppose the assumptions of \Cref{thm:main_convergence_hp} hold, and let
$W_U \in \R^{|\mathcal V| \times d}$ denote the unembedding matrix that
maps a $d$-dimensional latent vector to logits over the tokenizer
vocabulary, with operator norm $\norm{W_U}_{\mathrm{op}} \le B_U$. Define
the next-token entropy at the $\langle/\mathrm{think}\rangle$ position
after $T$ reasoning steps and the corresponding limit:
\begin{align*}
   H_T \;\coloneqq\; H\!\bigl(\softmax(W_U\, x_T)\bigr),
   \qquad
   H_\infty(Q) \;\coloneqq\; H\!\bigl(\softmax(W_U\, V^*(Q))\bigr),
\end{align*}
where $H(p) \coloneqq -\sum_i p_i \log p_i$ is the Shannon entropy of the
probability vector $p$. Then for every $Q \in F$ and every $T \ge 1$,
\begin{align}\label{eq:entropy_decay}
   \E\bigl[\, |H_T - H_\infty(Q)| \,\bigr]
   \;\le\;
   L_{\mathrm{sm}} \cdot B_U \cdot
   \E\!\bigl[\,\norm{x_T - V^*(Q)}\,\bigr],
\end{align}
where $L_{\mathrm{sm}}$ is the Lipschitz constant of the map
$\mathbf z \mapsto H(\softmax(\mathbf z))$ on the ball of radius
$B_U (M + \max_{Q \in F} \norm{V^*(Q)})$ in $\R^{|\mathcal V|}$
(Euclidean norm); $L_{\mathrm{sm}}$ is finite by the bounded-input
Lipschitz property of softmax-composed-with-entropy
(see \texttt{softmax-lipschitz.md}). Consequently
$|H_T - H_\infty(Q)|$ is bounded in expectation by an exponentially
decaying envelope in $T$:
\begin{align}\label{eq:entropy_decay_rate}
   \E\bigl[\, |H_T - H_\infty(Q)| \,\bigr]
   \;\le\;
   L_{\mathrm{sm}}\, B_U \cdot
   \Bigl(\tfrac{\gamma(Q)}{2} + \varepsilon_{\mathrm{anc}}\Bigr)
   \;+\;
   2 L_{\mathrm{sm}}\, B_U\, \bigl(M + \norm{V^*(Q)}\bigr)\,
   \exp\!\Bigl(-\frac{p_0 T}{8}\Bigr).
\end{align}
\end{corollary}

\begin{proof}
The proof is a four-line Lipschitz composition; see the technique digest
\texttt{softmax-lipschitz.md} for the underlying Lipschitz facts and
their references. (i) The softmax map
$\softmax: \R^{|\mathcal V|} \to \Delta^{|\mathcal V|}$ is $1$-Lipschitz
in $\ell_2 \to \ell_2$, because its Jacobian
$\mathrm{diag}(p) - p p^\top$ (with $p = \softmax(\mathbf z)$) has
operator norm at most $\max_i p_i (1 - p_i) \le 1/4 \le 1$.
(ii) The Shannon entropy $H: \Delta^{|\mathcal V|} \to \R$ is Lipschitz
on the image of softmax (which is bounded away from the simplex
boundary on any bounded set of inputs) with finite constant depending
on $|\mathcal V|$ and on the input bound. Let $L_{\mathrm{sm}}$ denote
the Lipschitz constant of the composition
$\mathbf z \mapsto H(\softmax(\mathbf z))$ on the relevant ball
(both $W_U x_T$ and $W_U V^*(Q)$ lie in a ball of radius
$B_U (M + \max_Q \norm{V^*(Q)})$ since $\norm{x_T} \le M$ by
\Cref{ass:bounded_value_norms} and the convex-combination argument of
\Cref{cor:expected_error_scaling}). (iii) The linear map
$\mathbf x \mapsto W_U \mathbf x$ has operator norm
$\norm{W_U}_{\mathrm{op}} \le B_U$, so composing:
$\norm{W_U x_T - W_U V^*(Q)} \le B_U \norm{x_T - V^*(Q)}$.
(iv) Therefore
$|H_T - H_\infty(Q)| = |H(\softmax(W_U x_T)) - H(\softmax(W_U V^*(Q)))|
\le L_{\mathrm{sm}} \cdot \norm{W_U x_T - W_U V^*(Q)}
\le L_{\mathrm{sm}} B_U \norm{x_T - V^*(Q)}$ pointwise. Taking
expectations gives \Cref{eq:entropy_decay}. Substituting
\Cref{eq:expected_error_scaling} into \Cref{eq:entropy_decay} yields
\Cref{eq:entropy_decay_rate}.
\end{proof}

\begin{remark}\label{rem:entropy_unembedding_norm}
The constant $B_U$ in \Cref{cor:entropy_decay} is the operator norm of
the unembedding matrix $W_U$, which is computable exactly from the
trained model's weights (one SVD computation). We treat it as a remark
rather than a numbered assumption because it is not load-bearing for
\Cref{thm:main_convergence_hp}; only \Cref{cor:entropy_decay} uses it.
The Lipschitz constant $L_{\mathrm{sm}}$ depends on $|\mathcal V|$ and
on the saturation level of the softmax (smaller when the softmax has
high-confidence components); for $\softmax \circ$ linear over a
finite vocabulary it is finite and computable for any given model.
\end{remark}

\begin{remark}\label{rem:choi2025_link}
\citet{choi2025entropy} document an empirical trajectory of $H_T$ ---
the next-token entropy at the $\langle/\mathrm{think}\rangle$ position
as a function of the reasoning length $T$ --- that decreases and
stabilises with thinking length, and they use the empirical plateau as
an early-exit signal. \Cref{cor:entropy_decay} predicts exactly this
trajectory: $H_T$ converges in expectation at exponential rate
$\exp(-p_0 T / 8)$ to a model-dependent floor
$H_\infty(Q) + \mathcal{O}(\gamma(Q)/2 + \varepsilon_{\mathrm{anc}})$.
The plateau threshold used in \cite{choi2025entropy} for early-exiting
is the empirical signature of this convergence; their EMA-based
stopping rule fires precisely when the exponentially-decaying envelope
of \Cref{eq:entropy_decay_rate} has shrunk to within the irreducible
floor.
\end{remark}

\begin{remark}[Test-time scaling law]\label{rem:test_time_scaling_law}
The exponential rate $\exp(-p_0 T / 8)$ in
\Cref{eq:test_time_scaling,eq:expected_error_scaling} is the test-time
analogue of the well-known training-time scaling laws: longer reasoning
trajectories systematically reduce the probability of an incorrect
answer, with the rate governed jointly by the anchor-emission
probability $p_0$ (the model's tendency to attend to answer-bearing
tokens) and the multiplicative-Chernoff exponent. The floor
$\gamma(Q)/2 + \varepsilon_{\mathrm{anc}}$ is the irreducible quality
budget set by the trained model: $\varepsilon_{\mathrm{anc}}$ controls
the anchor accuracy, and $\gamma(Q)/2$ controls how tightly the
trajectory has to land. The headline practical implication is that
test-time compute trades exponentially against accuracy, at a rate
determined entirely by inference-time-observable model quantities
($p_0$, $\Delta$, $\varepsilon_{\mathrm{anc}}$, $\gamma$).
\end{remark}

\begin{remark}[Per-layer/per-head extension]\label{rem:per_layer_extension}
\Cref{thm:main_convergence_hp} is stated for a single layer index $i$ and
head index $h$, with $V^*(Q)$ understood as the layer-$i$, head-$h$
target. To obtain a statement that holds simultaneously across all
$(i, h)$ pairs of a transformer with $L$ layers and $H$ heads per
layer (each pair has its own per-(layer, head) target $V^{*,(i,h)}(Q)$
and its own instantiation of
\Cref{ass:anchor_set_accuracy,ass:anchor_emission_prob,ass:score_margin,ass:bounded_value_norms,ass:decoding_existence}),
apply the union bound over the $L H$ per-(layer, head) failure events:
each pair's failure probability is at most $2 \exp(-p_0 T / 8)$, so
the joint failure probability is at most
$L H \cdot 2 \exp(-p_0 T / 8) = 2 L H \exp(-p_0 T / 8)$. To drive
the joint failure probability below a target $\delta \in (0, 1)$,
choose $T \ge (8/p_0) \log(2 L H / \delta)$ --- the standard
$\log(L H)$ multiplicative cost of multiplicity on top of the
single-(layer, head) horizon.
\end{remark}
